% Appendix H --- the similarities with the Loewner order. Moved out of section 5's "Smaller groups"
% whole (Q-0290, the author's decision of 2026-10-04; the review's length judgement, which keeps
% prop:similarity-cascade-families, the bound es <= 4 ecc and one summary paragraph in the main text):
% prop:similarity-stable-family, prop:similarity-orbit-family, rem:gaussian-matern-step,
% prop:similarity-general-sandwich, prop:similarity-general-bounded-shape,
% prop:similarity-degenerate-family, prop:similarity-line-family and
% prop:similarity-general-index-zero, with their `% shared with blueprint` markers, verbatim, and the
% sentences naming the cited facts they rest on. The numerical remark cites
% records/affine/check-gaussian-matern-step.py and its saved output; no statement uses a value.
% Q-0309 (fix group 2, review flag 37; records/affine/move-q0309-similarity-cascade.py):
% prop:similarity-cascade-families, every clause of which is about the statements here, moved from
% 5.7 to the end of this appendix, with the paragraph after it and its proof, verbatim.

\section{The similarities with the Loewner order}\label{app:similarity}

The statements of this appendix are summarized in \S\ref{sec:smaller-groups}, under ``The similarities
with the Loewner order'', and in the cell of Table~\ref{tab:landscape}. They concern families of
kernels from scale zero covariant under the similarities whose Loewner steps have infinitely
divisible increments, or positive ones when there is no Gaussian part. The question throughout is
how far from rotation invariance the non-Gaussian part of a kernel can be. The appendix answers it in
three steps. The first subsection takes the part to be elliptical, an image of one law, and bounds its
eccentricity. The second takes a general part, compares it with its limit at the thinnest apertures,
and shows which hypotheses bound its shape and which do not. The last subsection shows that every
family of the scheme has the infinitely divisible increments that the first two assume.

\subsection{An elliptical non-Gaussian part}

The first two statements take the non-Gaussian part to be elliptical.

% shared with blueprint prop:similarity-stable-family+297bca78b259
\begin{proposition}[similarity-covariant families with a stable part]\label{prop:similarity-stable-family}
Let $d \ge 2$ and $0 < \alpha < 2$. For maps $h : \PD(d) \to \PD(d)$ and $C : \PD(d) \to \PSD(d)$
put $\mu_\Sigma := N(0,C(\Sigma)) * S_\alpha(h(\Sigma))$, with $N(0,C)$ the centred Gaussian law of
covariance $C$.
\begin{enumerate}
\item \emph{(The stable parts are ordered.)} If $\Sigma, \Sigma' \in \PD(d)$ and
$\mu_{\Sigma'} = \mu_\Sigma * \rho$ for a positive finite measure $\rho$, then
$h(\Sigma) \preceq_\alpha h(\Sigma')$.
\item \emph{(Bounded shape.)} Suppose that $h(cQ\Sigma Q^{\tp}) = c\,Q\,h(\Sigma)\,Q^{\tp}$ for $c > 0$
and $Q \in \Orth(d)$, that $h$ extends continuously to $\PSD(d) \setminus \{0\}$ with values in
$\PD(d)$, and that every Loewner step $\Sigma \le \Sigma'$ in $\PD(d)$ has an increment as in (1).
Then $h(I) = h_0I$ with $h_0 > 0$, and with $c_1 := \lambda_{\max}(h(e_1e_1^{\tp})) > 0$,
\[
  \frac{\lambda_{\max}(h(\Sigma))}{\lambda_{\min}(h(\Sigma))} \le \Bigl(\frac{h_0}{c_1}\Bigr)^{\alpha}
  \qquad \text{for every } \Sigma \in \PD(d) .
\]
\item \emph{(A stable part that is not rotation invariant.)} Let $b \ge 0$ and $\kappa > 0$ with
$\alpha\kappa^2 - (d+1)\kappa - 1 \ge 0$, and put $h(\Sigma) := \Sigma + \kappa\tr(\Sigma)I$ and
$C(\Sigma) := 2b\Sigma$ for $\Sigma \in \PSD(d)$, with $S_\alpha(0) := \delta_0$. For every
$\Sigma \le \Sigma'$ in $\PSD(d)$ there is a probability measure $\rho_{\Sigma,\Sigma'}$ with
$\mu_{\Sigma'} = \mu_\Sigma * \rho_{\Sigma,\Sigma'}$, and the family $(\rho_{\Sigma,\Sigma'})$ satisfies
(K1), (K2) and (K4) of Definition~\ref{def:affine-cascade-family}, and (K3) for every similarity $A = cQ$,
$c > 0$, $Q \in \Orth(d)$. Its stable part $S_\alpha(h(\Sigma))$ is rotation invariant only when
$\Sigma$ is a multiple of $I$.
\end{enumerate}
\end{proposition}

In clause (3) the non-Gaussian part carries the shape of $\Sigma$, and $b = 0$ is allowed, so no
Gaussian part is needed for that. Clause (2)'s bound tends to $1$ as $\alpha \to 0$. For clause (3)'s family
that bound is $((1 + d\kappa)/(1 + \kappa))^\alpha$, and the eccentricity of the stable part is at most
$(1 + \kappa)/\kappa$, its value at rank one. The condition on
$\kappa$ is sufficient, not sharp. Without the extension of $h$ to rank one with values in $\PD(d)$
clause (2) fails (Proposition~\ref{prop:similarity-degenerate-family}).

% proof of prop:similarity-stable-family
\begin{proof}[Sketch]
\emph{(1).} The increment has the transform \ref{eq:stable-tail-transform} with
$\tfrac12(C(\Sigma') - C(\Sigma))$ for its matrix $C$, $H = h(\Sigma)$ and $H' = h(\Sigma')$, so
Lemma~\ref{lem:stable-tail} orders the \Levy{} densities, which is $h(\Sigma) \preceq_\alpha h(\Sigma')$ by
Proposition~\ref{prop:stable-step-cone}(1).

\emph{(2).} Equivariance under $\Orth(d)$ gives $h(I) = h_0I$, as in the proof of
Lemma~\ref{lem:affine-gauge}. Scale so that $\lambda_{\max}(\Sigma) = 1$, and sandwich
$Q\diag(1, \varepsilon, \dots, \varepsilon)Q^{\tp} \le \Sigma \le I$ with $Qe_1$ a top eigenvector. Let
$\eta_1 \le \dots \le \eta_d$ be the eigenvalues of $h(\Sigma)$. By (1) and the inclusion of
$\preceq_\alpha$ in the Loewner order, letting $\varepsilon \to 0$ gives
$Q\,h(e_1e_1^{\tp})\,Q^{\tp} \le h(\Sigma)$, so $\eta_d \ge c_1$. For $h(\Sigma) \preceq_\alpha h_0I$ the
relative step $M = h_0h(\Sigma)^{-1}$ has $\lambda_{\min}(M) = h_0/\eta_d$ and
$\det M = h_0^d/(\eta_1\cdots\eta_d)$, so condition (c) of Proposition~\ref{prop:stable-step-cone}(1) reads
\[
  \prod_i\frac{\eta_i}{\eta_d} \ge \Bigl(\frac{\eta_d}{h_0}\Bigr)^\alpha ,
  \qquad\text{whence}\qquad
  \frac{\eta_d}{\eta_1} \le \Bigl(\frac{h_0}{\eta_d}\Bigr)^\alpha \le \Bigl(\frac{h_0}{c_1}\Bigr)^\alpha ,
\]
the left-hand product being at most $\eta_1/\eta_d$.

\emph{(3).} Let $P := \Sigma' - \Sigma$ with $\Sigma \ne 0$, and put $H := h(\Sigma)$,
$\tau := \tr\Sigma$ and $u := \tr(P)/\tau$. The relative step
$M = H^{-1/2}h(\Sigma')H^{-1/2} = I + H^{-1/2}PH^{-1/2} + \kappa\tr(P)H^{-1}$ has
$\lambda_{\min}(M) \ge 1 + \kappa u/(1+\kappa)$ and $\tr M \le d + (1+d\kappa)u/\kappa$, because
$\kappa\tau I \le H \le (1+\kappa)\tau I$. With $A := (1+d\kappa)/\kappa$ and $B := \kappa/(1+\kappa)$
these bounds read $\tr M \le d + Au$ and $\lambda_{\min}(M) \ge 1 + Bu$, and the condition on $\kappa$,
$(d+\alpha)\kappa^2 \ge (1+d\kappa)(1+\kappa)$, is $A \le (d+\alpha)B$. The inequality of arithmetic and
geometric means for the eigenvalues of $M$, and Bernoulli's inequality, give
\[
  \det M \le \Bigl(\frac{\tr M}{d}\Bigr)^d
  \le \Bigl(1 + \frac{d+\alpha}{d}\,\frac{\kappa u}{1+\kappa}\Bigr)^d
  \le \lambda_{\min}(M)^{d+\alpha} ,
\]
which is condition (c). So the densities are ordered, and
Proposition~\ref{prop:stable-step-cone}(1) gives the stable increment; the Gaussian increment is
$N(0, 2bP)$. The axioms are read off the telescoping transforms.
\end{proof}

For an elliptical part of any self-decomposable profile, infinitely divisible increments order the
parts, and the bound on the shape is read off the dilation ratio.

% shared with blueprint prop:similarity-orbit-family+550d5a069682
\begin{proposition}[similarity-covariant families with an elliptical part: bounded shape, and shape-blind at index $0$]\label{prop:similarity-orbit-family}
Let $d \ge 2$, let $k$ be as in Proposition~\ref{prop:orbit-step-cone} and right-continuous, with dilation ratio
$R$, and write $\nu_H := (H^{1/2})_*\nu_I$, $H \in \PD(d)$, for the orbit of the law $\nu_I$ of
Proposition~\ref{prop:orbit-step-cone}(1), with \Levy{} densities $n_H$. For maps $h : \PD(d) \to \PD(d)$ and
$C : \PD(d) \to \PSD(d)$ put $\mu_\Sigma := N(0,C(\Sigma)) * \nu_{h(\Sigma)}$.
\begin{enumerate}
\item \emph{(Infinitely divisible increments order the parts.)} If $\Sigma, \Sigma' \in \PD(d)$ and
$\mu_{\Sigma'} = \mu_\Sigma * \rho$ with $\rho$ infinitely divisible, then $C(\Sigma) \le C(\Sigma')$
and $n_{h(\Sigma)} \le n_{h(\Sigma')}$ everywhere: the step $h(\Sigma) \to h(\Sigma')$ is in the step
set of Proposition~\ref{prop:orbit-step-cone}(2).
\item \emph{(Bounded shape.)} Suppose $h(cQ\Sigma Q^{\tp}) = c\,Q\,h(\Sigma)\,Q^{\tp}$ for $c > 0$ and
$Q \in \Orth(d)$, that $h$ extends continuously to $\PSD(d) \setminus \{0\}$ with values in $\PD(d)$,
and that every Loewner step $\Sigma \le \Sigma'$ in $\PD(d)$ has an infinitely divisible increment.
Then $h(I) = h_0I$ with $h_0 > 0$, and with $c_1 := \lambda_{\max}(h(e_1e_1^{\tp}))$,
\[
  \frac{\lambda_{\max}(h(\Sigma))}{\lambda_{\min}(h(\Sigma))} \le R\bigl((h_0/c_1)^{1/2}\bigr)^2
  \qquad \text{for every } \Sigma \in \PD(d) .
\]
\item \emph{(Index $0$: shape-blind.)} If $k(0{+})$ is finite, $h$ is equivariant as in (2) (no
extension needed), and every Loewner step has an infinitely divisible increment, then $h(\Sigma)$
is a multiple of $I$ for every $\Sigma$.
\item \emph{(Logarithmic exponent, no Gaussian part.)} Let instead the non-Gaussian parts be the
orbit $(\nu_H)$ of a law satisfying the hypothesis of Proposition~\ref{prop:log-orbit-rays}, let $C \equiv 0$
and $h$ be equivariant as in (2), and suppose every Loewner step has a positive finite increment.
Then $h(\Sigma)$ is a multiple of $I$ for every $\Sigma$.
\end{enumerate}
\end{proposition}

Without a Gaussian part, clause (4) gets shape-blindness from positive increments alone. With a
Gaussian part a positive increment need not order the parts, so clauses (1)--(3) require infinitely
divisible increments. A computation, Remark~\ref{rem:gaussian-matern-step}, suggests that this
hypothesis cannot simply be dropped; it is not a proof. Proposition~\ref{prop:similarity-cascade-families}
discharges the hypothesis for every family of the scheme.

% proof of prop:similarity-orbit-family
\begin{proof}[Sketch]
\emph{(1).} Both kernels are infinitely divisible with triplets $(C(\cdot), n_{h(\cdot)}(x)\,dx, 0)$, and
by the uniqueness C2(ii) the triplet of $\rho$ adds to that of $\mu_\Sigma$; so $C(\Sigma) \le C(\Sigma')$
and $n_{h(\Sigma)} \le n_{h(\Sigma')}$ almost everywhere. In the coordinates $x = h(\Sigma)^{1/2}y$, an
open set of positive measure near each radius and the right-continuity of $k$ give
$D\,m(r) \le m(r/\lambda)$ for every $r > 0$, where $m(r) = k(r)r^{-d}$, $D^2$ is the determinant of the
relative step $M$ and $\lambda^2$ its least eigenvalue; as $\langle y, M^{-1}y\rangle^{1/2} \le |y|/\lambda$
and $m$ is nonincreasing, that is the ordering everywhere.

\emph{(2)} is the proof of Proposition~\ref{prop:similarity-stable-family}(2), with the step set of
Proposition~\ref{prop:orbit-step-cone}(2) in place of the stable cone, $R$ being nondecreasing, and
$h(I) = h_0I$ as in the proof of Lemma~\ref{lem:affine-gauge}.

\emph{(3).} $R \equiv 1$ (Proposition~\ref{prop:orbit-step-cone}(3)), so every Loewner step moves $h$ along
a ray, and $\Sigma \le tI$ makes $h(\Sigma)$ a multiple of $h(tI) = th_0I$.

\emph{(4).} The positive increment of $\Sigma \le tI$ is an increment of the orbit $(\nu_H)$ from $h(\Sigma)$
to $th_0I$, and Proposition~\ref{prop:log-orbit-rays} makes them multiples.
\end{proof}

\begin{remark}[a Gaussian part and a positive increment, measured]\label{rem:gaussian-matern-step}
In the plane take $C(\Sigma) = 2b\Sigma$, the \Matern{} orbit of index $1$ as the elliptical part, and
$h(\Sigma) = \Sigma + \kappa\tr(\Sigma)I$ with $\kappa = 2$. Consider the step from $\diag(1, 0.3)$ by
$\diag(0.05, 0.55)$, for which $h(\Sigma')$ is not a multiple of $h(\Sigma)$. A script deposited
with the paper computes its would-be increment by the fast Fourier transform. It samples the transform
of the increment on the frequency grid of a spatial square of side $60\,\lambda_{\max}(h(\Sigma'))^{1/2}
\approx 132$ centred at the origin, at two resolutions, $1024^2$ and $2048^2$ points (spacings $0.129$
and $0.065$, frequencies up to $\pi$ over the spacing, about $24$ and $49$ in each coordinate). At
$b = 2$ the computed density is nonnegative to rounding:
its minimum is $-10^{-17}$ of its maximum at both resolutions, as for the ray steps that calibrate the
computation. At $b = 0.1$ it is negative, with negative mass $1.4\cdot10^{-2}$ at $2048^2$ points.
% Round 3 (content review), E item 8 (remark_h3_checks-output.txt): the value at the finer grid, and why.
At $1024^2$ points the Gaussian factor is still $0.05$ at the largest frequency. These are numerical
findings, and positivity at $b = 2$ is not proved. If that increment is positive, it is not infinitely
divisible, since its \Levy{} measure would be the signed $n_{h(\Sigma')} - n_{h(\Sigma)}$. The
computation thus suggests that the Gaussian part smooths away the obstruction of
Proposition~\ref{prop:log-orbit-rays} on this full-rank step. By
Proposition~\ref{prop:similarity-cascade-families} these kernels are the kernels from scale zero of no
(similarities, Loewner)-cascade family, so the example shows why the axioms require continuity along
the order.
\end{remark}

\subsection{A general non-Gaussian part}

For a non-Gaussian part with no orbit structure the same sandwich works in the order of \Levy{}
measures, and in the order of exponents when there is no Gaussian part.

% shared with blueprint prop:similarity-general-sandwich+86d7dd25f45f
\begin{proposition}[the sandwich for a general similarity-covariant part]\label{prop:similarity-general-sandwich}
Let $d \ge 2$. For each $\Sigma \in \PD(d)$ let $\mu_\Sigma$ be a symmetric infinitely divisible
probability measure on $\RR^d$ with generating triplet $(C(\Sigma), L_\Sigma, 0)$, and put
$\psi_\Sigma(\omega) := \int(1 - \cos\langle\omega,x\rangle)\,L_\Sigma(dx)$; so
$\mu_\Sigma = N(0,C(\Sigma)) * \nu_\Sigma$, where the \emph{non-Gaussian part} $\nu_\Sigma$ has
$\widehat\nu_\Sigma = e^{-\psi_\Sigma}$. Suppose the family is covariant under similarities,
$(\sqrt c\,Q)_*\mu_\Sigma = \mu_{cQ\Sigma Q^{\tp}}$ for $c > 0$, $Q \in \Orth(d)$ and
$\Sigma \in \PD(d)$, which is (K3) of Definition~\ref{def:affine-cascade-family} for $A = \sqrt c\,Q$. Put
$\Sigma_\varepsilon := \diag(1,\varepsilon,\dots,\varepsilon)$ for $0 < \varepsilon \le 1$.
\begin{enumerate}
\item \emph{(Covariance.)} $C(cQ\Sigma Q^{\tp}) = c\,Q\,C(\Sigma)\,Q^{\tp}$,
$L_{cQ\Sigma Q^{\tp}} = (\sqrt c\,Q)_*L_\Sigma$ and
$\psi_{cQ\Sigma Q^{\tp}}(\omega) = \psi_\Sigma(\sqrt c\,Q^{\tp}\omega)$. In particular $L_I$ and
$\psi_I$ are rotation invariant.
\item \emph{(Infinitely divisible increments.)} If $\mu_{\Sigma'} = \mu_\Sigma * \rho$ with $\rho$
infinitely divisible, then $C(\Sigma) \le C(\Sigma')$, $L_\Sigma \le L_{\Sigma'}$ and
$\psi_\Sigma \le \psi_{\Sigma'}$.
\item \emph{(Positive increments without a Gaussian part.)} If $C \equiv 0$ and
$\mu_{\Sigma'} = \mu_\Sigma * \rho$ with $\rho$ a positive finite measure, then $\rho$ is a
probability measure, $\widehat\rho = e^{-(\psi_{\Sigma'} - \psi_\Sigma)}$, and
$\psi_\Sigma \le \psi_{\Sigma'}$.
\item \emph{(The rank-one limit.)} Suppose every Loewner step $\Sigma \le \Sigma'$ in $\PD(d)$ has an
infinitely divisible increment. Then $L_*(B) := \lim_{\varepsilon\downarrow0}L_{\Sigma_\varepsilon}(B)$
exists for every Borel set $B$ bounded away from the origin, and $L_*$ extends to a \Levy{} measure
with $L_* \le L_{\Sigma_\varepsilon}$ for every $\varepsilon$; and
$\psi_*(\omega) := \lim_{\varepsilon\downarrow0}\psi_{\Sigma_\varepsilon}(\omega)$ exists for every
$\omega$. Both are invariant under every $R \in \Orth(d)$ with $Re_1 = \pm e_1$. If instead
$C \equiv 0$ and every Loewner step has a positive finite increment, $\psi_*$ exists and has the same
invariance.
\item \emph{(The sandwich.)} Let $\Sigma \in \PD(d)$ with $\lambda_{\max}(\Sigma) = 1$, and let
$Q \in \Orth(d)$ have as its first column a unit eigenvector of $\Sigma$ for the eigenvalue $1$. Under
the first hypothesis of (4),
\[
  Q_*L_* \le L_\Sigma \le L_I
  \qquad\text{and}\qquad
  \psi_*(Q^{\tp}\omega) \le \psi_\Sigma(\omega) \le \psi_I(\omega) \quad (\omega \in \RR^d) ;
\]
under the second, the inequalities between the exponents hold.
\end{enumerate}
\end{proposition}

The rank-one limit needs no continuity hypothesis: the \Levy{} measures along $\Sigma_\varepsilon$
decrease as $\varepsilon \downarrow 0$, so the limit exists by monotonicity alone. The analogue of the
extension of $h$ to rank one in the elliptical statements is therefore a hypothesis on the size of the
limit, and the limit may be $0$ (Proposition~\ref{prop:similarity-degenerate-family}). For the families of the scheme the standing hypotheses
and the first hypothesis of (4) are theorems (Proposition~\ref{prop:similarity-cascade-families}(1)).

% proof of prop:similarity-general-sandwich
\begin{proof}[Sketch]
\emph{(1).} The symmetry of $\mu_\Sigma$ and the uniqueness C2(ii) make $L_\Sigma$ symmetric, so the
transform collapses to $\exp[-\tfrac12\langle\omega,C(\Sigma)\omega\rangle - \psi_\Sigma(\omega)]$ as in the
proof of Lemma~\ref{lem:elliptical-stable}(2). The image under $A = \sqrt c\,Q$, which (K3) of
Definition~\ref{def:affine-cascade-family} identifies with $\mu_{cQ\Sigma Q^{\tp}}$, has the triplet
$(AC(\Sigma)A^{\tp}, A_*L_\Sigma, 0)$, and uniqueness gives the identities.

\emph{(2).} The triplets of a convolution add, as in the proof of
Proposition~\ref{prop:similarity-orbit-family}(1). \emph{(3).} The transform of $\mu_\Sigma$ is positive, so
$\widehat\rho = e^{-(\psi_{\Sigma'} - \psi_\Sigma)}$ and $|\widehat\rho| \le 1$.

\emph{(4).} $\varepsilon \mapsto L_{\Sigma_\varepsilon}(B)$ is nondecreasing and bounded by $L_I(B)$ on sets
bounded away from the origin, so it converges. On a shell $\{|x| \ge \delta\}$ the limit is countably
additive, since for disjoint $B_n$ the terms $L_{\Sigma_\varepsilon}(B_n)$ are dominated by the summable
$L_I(B_n)$, and it extends to a \Levy{} measure below every $L_{\Sigma_\varepsilon}$. The invariance holds because $R\Sigma_\varepsilon R^{\tp} =
\Sigma_\varepsilon$ when $Re_1 = \pm e_1$.

\emph{(5).} $Q\Sigma_\varepsilon Q^{\tp} \le \Sigma \le I$ for $\varepsilon \le \lambda_{\min}(\Sigma)$, and (2) or
(3) with (1) and (4) gives the inequalities.
\end{proof}

The shape of a general part is measured by one of three functionals, and each bound is the sandwich
read in one of them.

% shared with blueprint prop:similarity-general-bounded-shape+892ef79455a8
\begin{proposition}[bounded shape under a comparable rank-one limit]\label{prop:similarity-general-bounded-shape}
For a measure $L$ on $\RR^d \setminus \{0\}$, a function $\psi \ge 0$ on $\RR^d$, $s > 0$ and
$H \in \PD(d)$, put
\[
  \operatorname{rs}(L) := \inf\{K \ge 1 : Q_*L \le K\,L \text{ for every } Q \in \Orth(d)\},
\]
\[
  \operatorname{es}(\psi) := \inf\{K \ge 1 : \psi(Q\omega) \le K\,\psi(\omega) \text{ for every } Q \in \Orth(d),\ \omega \in \RR^d\},
\]
\[
  \mathcal M_s(L) := \int(|x|^2 \wedge s^2)\,\frac{xx^{\tp}}{|x|^2}\,L(dx),
  \qquad T_s(L) := \tr\mathcal M_s(L) = \int(|x|^2 \wedge s^2)\,L(dx),
\]
\[
  \operatorname{ecc}(H) := \frac{\lambda_{\max}(H)}{\lambda_{\min}(H)} ,
\]
the \emph{rotational spread}, the \emph{exponent spread}, the \emph{moment matrix at scale $s$} and
the eccentricity, the infimum of the empty set being $\infty$, and $\operatorname{cond}M := \lambda_{\max}(M)/\lambda_{\min}(M)
\in [1,\infty]$ for $M \in \PSD(d) \setminus \{0\}$. Let the family be as in
Proposition~\ref{prop:similarity-general-sandwich}.
\begin{enumerate}
\item \emph{(The functionals.)} For $c > 0$ and $Q \in \Orth(d)$,
$\operatorname{rs}((cQ)_*L) = \operatorname{rs}(L)$, $\operatorname{es}(\psi(cQ\,\cdot)) = \operatorname{es}(\psi)$
and $\mathcal M_s((cQ)_*L) = c^2Q\,\mathcal M_{s/c}(L)\,Q^{\tp}$. $\operatorname{rs}(L) = 1$ exactly when $L$ is rotation
invariant, $\operatorname{es}(\psi) = 1$ exactly when $\psi$ is, and $\operatorname{es}(\psi_L) \le \operatorname{rs}(L)$ for the
exponent $\psi_L(\omega) := \int(1 - \cos\langle\omega,x\rangle)\,L(dx)$; the exponent spread of a sum is at
most the largest exponent spread of its terms. For the orbit $(\nu_H)$ of a profile as in
Proposition~\ref{prop:orbit-step-cone}(1), $\operatorname{es}(\psi_{\nu_H}) \le 4\operatorname{ecc}(H)$; for the stable orbit,
$\operatorname{es}(\psi_{S_\alpha(H)}) = \operatorname{ecc}(H)^{\alpha/2}$.
\item \emph{(\Levy{} level.)} Suppose every Loewner step in $\PD(d)$ has an infinitely divisible
increment and $L_* \ge \theta L_I$ for some $\theta \in (0,1]$. Then for every $\Sigma \in \PD(d)$, with
$t := \lambda_{\max}(\Sigma)$,
\[
  \theta\,L_{tI} \le L_\Sigma \le L_{tI} ,
  \qquad
  \operatorname{rs}(L_\Sigma) \le 1/\theta ,
  \qquad
  \operatorname{es}(\psi_\Sigma) \le 1/\theta .
\]
\item \emph{(Exponent level.)} Suppose every Loewner step in $\PD(d)$ has an infinitely divisible
increment, or that $C \equiv 0$ and every Loewner step has a positive finite increment; and that
$\psi_* \ge \theta\psi_I$ for some $\theta \in (0,1]$. Then $\theta\psi_{tI} \le \psi_\Sigma \le \psi_{tI}$ and
$\operatorname{es}(\psi_\Sigma) \le 1/\theta$ for every $\Sigma \in \PD(d)$, $t := \lambda_{\max}(\Sigma)$.
\item \emph{(Moments, scale by scale.)} Suppose every Loewner step in $\PD(d)$ has an infinitely
divisible increment, and $L_I \ne 0$. For $s > 0$ the limit $\mathcal M_s^* := \lim_{\varepsilon\downarrow0}\mathcal M_s(L_{\Sigma_\varepsilon})$
exists and is $\diag(a_s, b_s, \dots, b_s)$; put $\theta_s := d\min(a_s,b_s)/T_s(L_I) \in [0,1]$. Then for every
$\Sigma \in \PD(d)$ and $s > 0$, with $t := \lambda_{\max}(\Sigma)$,
$\operatorname{cond}\mathcal M_s(L_\Sigma) \le 1/\theta_{s/\sqrt t}$; so
$\operatorname{cond}\mathcal M_s(L_\Sigma) \le 1/\inf_{r>0}\theta_r$ for all $\Sigma$ and $s$.
\end{enumerate}
\end{proposition}

The exponent spread is the functional in which the shape is measured. It is defined for every
symmetric infinitely divisible law, it is the condition number for a Gaussian, it is
$\operatorname{ecc}(H)^{\alpha/2}$ on the stable orbit and at most $4\operatorname{ecc}(H)$ on the
elliptical orbit of every nonincreasing profile, so it extends the eccentricity that the elliptical statements bound, and it equals
$1$ exactly for the rotation-invariant parts. It is also the only one of the three that positive
increments see when there is no Gaussian part (clause (3)). The hypotheses are sufficient, not
necessary: the shape-blind family $((\det\Sigma)^{1/(2d)})_*\nu_I$ has $L_* = 0$ and spread $1$.

% proof of prop:similarity-general-bounded-shape
\begin{proof}[Sketch]
\emph{(1).} The rotational and the exponent spread are unchanged under conjugation and dilation, and
$\mathcal M_s((cQ)_*L) = c^2Q\,\mathcal M_{s/c}(L)\,Q^{\tp}$ by the substitution $x \mapsto cQx$. If
$Q_*L \le L$ for every $Q$, then $L = Q_*(Q^{-1})_*L \le Q_*L$, so equality holds. For the orbit of a profile, $\psi_{\nu_H}(\omega) = g(\langle\omega,H\omega\rangle)$
with $g$ nondecreasing, as $k$ is nonincreasing in the substitution of
Proposition~\ref{prop:orbit-step-cone}(1), and $g(4u) \le 4g(u)$ from $1 - \cos2\tau \le 4(1 - \cos\tau)$;
so $g(cu) \le 4c\,g(u)$ for $c \ge 1$, which gives $4\operatorname{ecc}(H)$. For the stable orbit the
exponent is $\langle\omega,H\omega\rangle^{\alpha/2}$ (Lemma~\ref{lem:elliptical-stable}(2)).

\emph{(2), (3).} After scaling to $t = 1$, Proposition~\ref{prop:similarity-general-sandwich}(5) gives
$\theta L_I \le Q_*L_* \le L_\Sigma \le L_I$, and, for every $Q' \in \Orth(d)$,
$Q'_*L_\Sigma \le L_I \le \theta^{-1}L_\Sigma$; the same for exponents.

\emph{(4).} $\mathcal M_s$ is monotone in the \Levy{} measure, so the limit exists; its invariance under
the $Q'$ with $Q'e_1 = \pm e_1$ makes it $\diag(a_s, b_s, \dots, b_s)$, by the argument of the proof of
Lemma~\ref{lem:affine-gauge}, and $\mathcal M_s(L_I) = \tfrac1dT_s(L_I)\,I$. The sandwich gives
$Q\mathcal M_s^*Q^{\tp} \le \mathcal M_s(L_\Sigma) \le \tfrac1dT_s(L_I)\,I$.
\end{proof}

Without a hypothesis on the rank-one limit the shape need not be bounded, already for elliptical
stable parts.

% shared with blueprint prop:similarity-degenerate-family+82099345e828
\begin{proposition}[without a nondegenerate rank-one limit the shape is unbounded]\label{prop:similarity-degenerate-family}
Let $d \ge 2$, $0 < \alpha < 2$, $\tfrac{1}{d+\alpha} \le \gamma < \tfrac1d$ and $p := 1 - d\gamma$, so that
$0 < p \le \tfrac{\alpha}{d+\alpha}$. For $\Sigma \in \PSD(d)$ put $h(\Sigma) := (\det\Sigma)^\gamma\,\Sigma^p$,
which is $0$ when $\Sigma$ is singular. For $b \ge 0$ put
$\mu_\Sigma := N(0,2b\Sigma) * S_\alpha(h(\Sigma))$, with $S_\alpha$ and $n_H$ as in
Lemma~\ref{lem:elliptical-stable} and $S_\alpha(0) := \delta_0$; write $L_\Sigma := n_{h(\Sigma)}(x)\,dx$ and
$\psi_\Sigma$ for the \Levy{} measure and the exponent of the non-Gaussian part, $\Sigma \in \PD(d)$, and
use the notation of Proposition~\ref{prop:similarity-general-bounded-shape}.
\begin{enumerate}
\item \emph{(A cascade.)} $h$ is continuous on $\PSD(d)$, $h(cQ\Sigma Q^{\tp}) = c\,Q\,h(\Sigma)\,Q^{\tp}$ for
$c > 0$ and $Q \in \Orth(d)$, and $n_{h(\Sigma)} \le n_{h(\Sigma')}$ for every $\Sigma \le \Sigma'$ in $\PD(d)$.
For every $\Sigma \le \Sigma'$ in $\PSD(d)$ there is an infinitely divisible probability measure
$\rho_{\Sigma,\Sigma'}$ with $\mu_{\Sigma'} = \mu_\Sigma * \rho_{\Sigma,\Sigma'}$, and the family
$(\rho_{\Sigma,\Sigma'})$ satisfies (K1), (K2) and (K4) of Definition~\ref{def:affine-cascade-family}, and (K3) for
every similarity $A = cQ$.
\item \emph{(Unbounded shape.)} For $\Sigma \in \PD(d)$, $\operatorname{ecc}(h(\Sigma)) = \operatorname{ecc}(\Sigma)^p$,
$\operatorname{es}(\psi_\Sigma) = \operatorname{ecc}(\Sigma)^{p\alpha/2}$, $\operatorname{rs}(L_\Sigma) \ge \operatorname{ecc}(\Sigma)^{p\alpha/2}$, and for
every $s > 0$
\[
  \operatorname{cond}\mathcal M_s(L_\Sigma) \ge \frac{G_1}{d\,G_0}\operatorname{ecc}(\Sigma)^{p\alpha/2} ,
\]
\[
  G_0 := \int_{\RR^d}(z_d^2 \wedge 1)\,|z|^{-d-\alpha}\,dz ,
  \quad
  G_1 := \int_{\RR^d}(|z|^2 \wedge 1)\,|z|^{-d-\alpha}\,dz ,
\]
both finite and positive. The rank-one limit of Proposition~\ref{prop:similarity-general-sandwich}(4) is $L_* = 0$.
\item \emph{(A nonzero rotation-invariant rank-one limit is not enough.)} Let $0 < \beta < \alpha$, let
$S'_\beta(H)$ and $n'_H$ be the law and the \Levy{} density of Lemma~\ref{lem:elliptical-stable} for the index
$\beta$, with $S'_\beta(0) := \delta_0$, and put $\mu'_\Sigma := \mu_\Sigma * S'_\beta(\tr(\Sigma)\,I)$, with \Levy{} measure
$L'_\Sigma = (n_{h(\Sigma)} + n'_{\tr(\Sigma)I})(x)\,dx$ and exponent $\psi'_\Sigma$. The family $(\mu'_\Sigma)$
has all the properties of (1); its rank-one limit is $L'_* = n'_I(x)\,dx$, nonzero and rotation
invariant; for every $s > 0$, $\theta_s > 0$ in Proposition~\ref{prop:similarity-general-bounded-shape}(4); and for
every $\Sigma \in \PD(d)$
\[
  \operatorname{es}(\psi'_\Sigma) \ge \operatorname{ecc}(\Sigma)^{p\alpha/2} ,
  \qquad
  \sup_{s>0}\operatorname{cond}\mathcal M_s(L'_\Sigma) \ge \frac{G_1}{d\,G_0}\operatorname{ecc}(\Sigma)^{p\alpha/2} .
\]
\end{enumerate}
\end{proposition}

So under similarity covariance with a cascade along the whole Loewner order, the non-Gaussian part
need not stay within a bounded neighbourhood of isotropy, with or without a Gaussian part. The
extension of $h$ to rank one with values in $\PD(d)$ cannot be dropped: here $h$ extends by $0$ and the
shape is unbounded. It is not necessary either, since the gauge $(\det\Sigma)^{1/d}I$, the endpoint
$\gamma = 1/d$, also vanishes at rank one and has spread $1$. The family interpolates between the stable
orbit, $\gamma = 0$, which has no rank-one step, and that shape-blind gauge. Clause (3) shows that a
nonzero rank-one limit, even a rotation-invariant one, is not enough: what
Proposition~\ref{prop:similarity-general-bounded-shape} needs is comparability with $L_I$.

% proof of prop:similarity-degenerate-family
\begin{proof}[Sketch]
\emph{The Loewner--Heinz inequality} (\sscite{bhatia1997matrix}, Thm.~V.1.9, p.~115). For $0 < p < 1$,
$\Sigma^p = \kappa_p^{-1}\int_0^\infty(I - s(\Sigma + sI)^{-1})s^{p-1}\,ds$, and inversion reverses the
order on $\PD(d)$, so $\Sigma \le \Sigma'$ gives $\Sigma^p \le \Sigma'^p$ and $\Sigma'^{-p} \le \Sigma^{-p}$.

\emph{(1).} $\det h(\Sigma) = \det\Sigma$, so by \ref{eq:stable-levy-density}
\[
  n_{h(\Sigma)}(x) = c_{d,\alpha}^{-1}\,(\det\Sigma)^{((d+\alpha)\gamma - 1)/2}\,\langle x,\Sigma^{-p}x\rangle^{-(d+\alpha)/2} ,
\]
and along a Loewner step both factors increase, the first because $(d+\alpha)\gamma \ge 1$. So
Proposition~\ref{prop:stable-step-cone}(1) gives the stable increments, and the axioms follow as in the
proof of Proposition~\ref{prop:similarity-stable-family}(3). For $\gamma < 1/(d+\alpha)$ the display gives
$n_{h(I + \tau e_1e_1^{\tp})}(e_2) < n_{h(I)}(e_2)$, and Lemma~\ref{lem:stable-tail} excludes a positive
increment, so the lower bound on $\gamma$ is needed.

\emph{(2).} The eigenvalues of $h(\Sigma)$ are $(\det\Sigma)^\gamma\sigma_i^p$, $\sigma_i$ those of $\Sigma$,
so $\operatorname{ecc}(h(\Sigma)) = \operatorname{ecc}(\Sigma)^p$, and
Proposition~\ref{prop:similarity-general-bounded-shape}(1) gives the spreads. For the moments let $H \in \PD(d)$ have
unit eigenvectors $v_\pm$ for its eigenvalues $\eta_+ := \lambda_{\max}(H)$ and
$\eta_- := \lambda_{\min}(H)$, and substitute $x = H^{1/2}y$, under which
$n_H(x)\,dx = c_{d,\alpha}^{-1}|y|^{-d-\alpha}\,dy$ and $\langle v_\pm,x\rangle = \eta_\pm^{1/2}\langle v_\pm,y\rangle$.
A rotation taking $v_-$ to $e_d$ and the dilation $y = s\eta_-^{-1/2}z$ give the upper bound, and
$|x|^2 \le \eta_+|y|^2$ with the average $\tfrac1d$ of $\langle v_+,y\rangle^2/|y|^2$ over the sphere the lower:
\[
  \langle v_-,\mathcal M_s(n_H)v_-\rangle \le c_{d,\alpha}^{-1}\eta_-^{\alpha/2}s^{2-\alpha}G_0 ,
  \qquad
  \langle v_+,\mathcal M_s(n_H)v_+\rangle \ge \frac{c_{d,\alpha}^{-1}}{d}\,\eta_+^{\alpha/2}s^{2-\alpha}G_1 .
\]
Their ratio is at most $\operatorname{cond}\mathcal M_s(n_H)$, which gives the bound with $H = h(\Sigma)$.
Along $\Sigma_\varepsilon$ the densities tend to $0$ off a null set, so $L_* = 0$ by dominated convergence,
in the sense of Proposition~\ref{prop:similarity-general-sandwich}(4).

\emph{(3).} The gauge $\tr(\Sigma)I$ is equivariant and monotone, so the second factor cascades too, and
$L'_* = n'_I(x)\,dx$. At large $|\omega|$ the index-$\alpha$ term dominates the exponent, and at small $s$ it
dominates the moments, which gives the two bounds.
\end{proof}

The rotational spread is not equivalent to the other two functionals, and exact shape-blindness at index
$0$ fails for a general part. Both are shown by a part laid along lines.

% shared with blueprint prop:similarity-line-family+b0f25a81e985
\begin{proposition}[a part laid along lines]\label{prop:similarity-line-family}
Let $d \ge 2$, let $k : (0,\infty) \to [0,\infty)$ be nonincreasing and not identically $0$ with
$\int_0^1rk(r)\,dr + \int_1^\infty k(r)\,dr/r < \infty$, let $\ell(dr) := k(|r|)\,dr/|r|$ on
$\RR \setminus \{0\}$, and let $\overline\sigma$ be the normalised surface measure on $S^{d-1}$. For
$\Sigma \in \PSD(d)$ and $q_\Sigma(u) := \langle u,\Sigma u\rangle$ let $L_\Sigma$ be the measure on
$\RR^d \setminus \{0\}$ with
\[
  \int f\,dL_\Sigma := \int_{S^{d-1}}\overline\sigma(du)\int_{\RR}f\bigl(r\,q_\Sigma(u)^{1/2}u\bigr)\,\ell(dr)
\]
for Borel $f \ge 0$ on $\RR^d$ with $f(0) = 0$, and for $b \ge 0$ let $\mu_\Sigma := N(0,2b\Sigma) * \nu_\Sigma$,
$\nu_\Sigma$ the law with triplet $(0, L_\Sigma, 0)$.
\begin{enumerate}
\item \emph{(A cascade.)} $L_\Sigma$ is a symmetric \Levy{} measure; for $\Sigma \in \PD(d)$ it is the image
of $\overline\sigma(du)\,2k(s\,q_\Sigma(u)^{-1/2})\,ds/s$ under $(u,s) \mapsto su$, a polar profile nonincreasing in
$s$ and at most $2k(0{+})$; $L_{cQ\Sigma Q^{\tp}} = (\sqrt c\,Q)_*L_\Sigma$; and $L_\Sigma \le L_{\Sigma'}$ for
$\Sigma \le \Sigma'$. Hence for every $\Sigma \le \Sigma'$ in $\PSD(d)$ there is an infinitely divisible
probability measure $\rho_{\Sigma,\Sigma'}$ with $\mu_{\Sigma'} = \mu_\Sigma * \rho_{\Sigma,\Sigma'}$, and the family
$(\rho_{\Sigma,\Sigma'})$ satisfies (K1), (K2), (K4) of Definition~\ref{def:affine-cascade-family} and (K3) for every
similarity.
\item \emph{(Three functionals.)} For every $\Sigma \in \PD(d)$ and $s > 0$,
$\operatorname{es}(\psi_\Sigma) \le 4d(d+2)$ and $\operatorname{cond}\mathcal M_s(L_\Sigma) \le d+2$, while
$\operatorname{rs}(L_{\Sigma_\varepsilon}) \to \infty$ as $\varepsilon \downarrow 0$. The rank-one limit $L_*$ is at least the
image of $\overline\sigma(du)\,2k(s/|u_1|)\,ds/s$ (the set $u_1 = 0$ is $\overline\sigma$-null), so it
charges every open cone.
\item \emph{(Not shape-blind at index $0$.)} If $\Sigma \in \PD(d)$ is not a multiple of $I$, then $L_\Sigma$,
and so $\nu_\Sigma$, is not rotation invariant. This holds in particular when $k(0{+}) < \infty$, so that
every polar profile of every $L_\Sigma$ is bounded at the origin.
\end{enumerate}
\end{proposition}

The \Levy{} measure is laid along lines, and in direction $u$ it is the one-dimensional profile dilated
by $q_\Sigma(u)^{1/2}$. The exponent and the moments stay within a bounded neighbourhood of isotropy,
while the rotational spread does not, although the rank-one limit is nonzero and charges every
direction. Clause (3) holds even under infinitely divisible increments: a general part with polar
profiles bounded at the origin need not be shape-blind. For this family the index at the origin is
$2k(0{+})$ in every direction for every $\Sigma$, and the shape is carried at positive radii.

% proof of prop:similarity-line-family
\begin{proof}[Sketch]
\emph{(1).} Splitting the line integral at $0$ and substituting $s = rq^{1/2}$ gives the polar form,
nonincreasing in $s$ because $k$ is, and $q_\Sigma \le q_{\Sigma'}$ for $\Sigma \le \Sigma'$ orders the
measures. C2(iii) gives the laws and increments, and the axioms (K1)--(K4) of
Definition~\ref{def:affine-cascade-family}, (K3) for similarities, follow from the transforms as in the
proof of Proposition~\ref{prop:similarity-general-sandwich}(1).

% CHANGED Q-0281 wave 2: blind review W2I-4: the sketch of (2) re-transcribed to the route of the
% proof of record: the bounds at Sigma itself from the fourth moments, fixed caps for the rotational
% spread, open cones through sectors. It no longer routes through
% prop:similarity-general-bounded-shape(3)-(4), and reads no law.
\emph{(2).} The exponent is $\psi_\Sigma(\omega) = \int\overline\sigma(du)\,\psi_\ell(q_\Sigma(u)^{1/2}\langle\omega,u\rangle)$,
where $\psi_\ell$, the exponent of the line profile, is even, nondecreasing on $[0,\infty)$ and satisfies
$\psi_\ell(a\tau) \ge \tfrac{a^2}{4}\psi_\ell(\tau)$ for $0 < a \le 1$, from $\psi_\ell(2\tau) \le
4\psi_\ell(\tau)$ and the monotonicity of $\psi_\ell$. Let $\Sigma \in \PD(d)$, $t := \lambda_{\max}(\Sigma)$ and
$v$ a unit eigenvector for it, so that $t\langle v,u\rangle^2 \le q_\Sigma(u) \le t$ on the sphere. Then
$\psi_\Sigma(R\omega) \le \psi_\ell(t^{1/2}|\omega|)$ for every $R \in \Orth(d)$. For $\omega \ne 0$, $w := \omega/|\omega|$ and
$a := |\langle v,u\rangle|\,|\langle w,u\rangle| \le 1$, the fourth moments of the normalised surface measure,
$\int\langle v,u\rangle^2\langle w,u\rangle^2\,\overline\sigma(du) = \frac{1 + 2\langle v,w\rangle^2}{d(d+2)}$, give
\[
  \psi_\Sigma(\omega) \ge \int\overline\sigma(du)\,\psi_\ell\bigl(a\,t^{1/2}|\omega|\bigr)
  \ge \tfrac14\psi_\ell(t^{1/2}|\omega|)\int\langle v,u\rangle^2\langle w,u\rangle^2\,\overline\sigma(du)
  \ge \frac{\psi_\ell(t^{1/2}|\omega|)}{4d(d+2)} ,
\]
which is the bound on $\operatorname{es}(\psi_\Sigma)$, the functionals being those of
Proposition~\ref{prop:similarity-general-bounded-shape}. In the same way
$\langle y,\mathcal M_s(L_\Sigma)y\rangle = \int\overline\sigma(du)\,\langle y,u\rangle^2G_s(q_\Sigma(u))$, with
$G_s(\tau) := \int(\tau r^2 \wedge s^2)\,\ell(dr)$, lies between $G_s(t)|y|^2/(d(d+2))$ and $G_s(t)|y|^2/d$, so
$\operatorname{cond}\mathcal M_s(L_\Sigma) \le d+2$. For the rotational spread fix $r_0$ with $k(r_0) > 0$, the shell
$r_0/4 < |x| < r_0/2$ and, for $q_2 \in (0,\tfrac14]$, the cap $U$ of directions with $u_1^2 < q_2/2$ and
$u_2^2 > \tfrac14$; let $R$ exchange $e_1$ and $e_2$. For $\varepsilon < q_2/2$, $q_{\Sigma_\varepsilon} < q_2$ on $U$ and
$q_{\Sigma_\varepsilon} > \tfrac14$ on $RU$, so on the shell the sector over $U$ carries at most $2\overline\sigma(U)I(q_2)$ and
the sector over $RU$ at least $2\overline\sigma(U)I(\tfrac14) > 0$, where
$I(q) := \int_{r_0/4}^{r_0/2}k(sq^{-1/2})\,ds/s \le \log 2\cdot k(\tfrac{r_0}{4}q^{-1/2})$ tends to $0$ with $q$, $k$ tending
to $0$ at infinity. So the rotational spread is unbounded. Finally $q_{\Sigma_\varepsilon}(u) \ge u_1^2$ puts the image of
$\overline\sigma(du)\,2k(s/|u_1|)\,ds/s$ below every $L_{\Sigma_\varepsilon}$, hence below $L_*$; and a nonempty open cone
contains a sector $\{x : cr_0/2 < |x| < cr_0,\ x/|x| \in V\}$ with $V$ open, $\overline\sigma(V) > 0$ and $|u_1| > c$ on
$V$, which every $L_{\Sigma_\varepsilon}$ charges with mass at least $\overline\sigma(V)k(r_0)\log 2$.

\emph{(3).} Let $\lambda_{\min}(\Sigma) < m'' < m' < \lambda_{\max}(\Sigma)$. If $L_\Sigma$ were rotation
invariant, a cone of directions around the largest eigenvector, on which $q_\Sigma > m'$, and its
image under the reflection exchanging the largest and the least eigenvectors, on which
$q_\Sigma < m''$, would carry equal mass on every shell. As $k$ is nonincreasing,
this gives $\int_a^b k(s/\sqrt{m'})\,ds/s \le \int_a^b k(s/\sqrt{m''})\,ds/s$ for $0 < a < b$, the reverse of
the pointwise inequality between the integrands. So $k(r) = k(\lambda r)$ for almost every $r$, with
$\lambda := (m'/m'')^{1/2} > 1$, and then $k = 0$.
\end{proof}

What survives of shape-blindness at index $0$ for a general part is shape-blindness at the origin.

% shared with blueprint prop:similarity-general-index-zero+64089ae1c63d
\begin{proposition}[index $0$ for a general part: shape-blind at the origin]\label{prop:similarity-general-index-zero}
Let the family be as in Proposition~\ref{prop:similarity-general-sandwich}, and suppose $L_I$ is the \Levy{} measure of
an isotropic origin datum $(0, k_d)$ in dimension $d$ (Definition~\ref{def:isotropic-origin-datum}) with
$c := k_d(0{+}) < \infty$. Let $\overline\sigma$ be the normalised surface measure on $S^{d-1}$, and for a Borel
set $A \subseteq S^{d-1}$ and $0 < \varepsilon < 1$ put $K_A(\varepsilon) := \{x : \varepsilon < |x| \le 1,\ x/|x| \in A\}$.
\begin{enumerate}
\item \emph{(Infinitely divisible increments: the index at the origin.)} If every Loewner step in $\PD(d)$ has
an infinitely divisible increment, then for every $\Sigma \in \PD(d)$ and every Borel $A \subseteq S^{d-1}$
\[
  \lim_{\varepsilon\downarrow0}\frac{L_\Sigma(K_A(\varepsilon))}{\log(1/\varepsilon)} = c\,\overline\sigma(A) .
\]
\item \emph{(The exponent, to within a bounded term.)} If every Loewner step has an infinitely divisible
increment, or if $C \equiv 0$ and every Loewner step has a positive finite increment, then for every
$\Sigma \in \PD(d)$ every limit point of $\psi_\Sigma(\omega) - \psi_I(\omega)$ as $|\omega| \to \infty$ lies in
$[\tfrac c2\log\lambda_{\min}(\Sigma), \tfrac c2\log\lambda_{\max}(\Sigma)]$.
\item \emph{(The blow-up.)} Under either hypothesis of (2), let $\Sigma \in \PD(d)$ and suppose
\emph{(D)}: for $\overline\sigma$-almost every $u \in S^{d-1}$ the limit
$\Delta_\Sigma(u) := \lim_{r\to\infty}\bigl(\psi_\Sigma(ru) - \psi_I(ru)\bigr)$ exists. Then $\Delta_\Sigma$ is
$\overline\sigma$-almost everywhere equal to a constant.
\end{enumerate}
\end{proposition}

Under infinitely divisible increments, with or without a Gaussian part, the \Levy{} measure of every
part has the isotropic member's logarithmic index $c$ at the origin, spread uniformly over directions
(clause (1)). Under positive increments alone, and without a Gaussian part, the exponents are
sandwiched in the same way (clause (2)), and the blow-up of Proposition~\ref{prop:log-orbit-rays} works
for a general part under (D) (clause (3)). Condition (D) for a general part with polar profiles bounded
at the origin requires the second term of the exponent's expansion, which needs more than boundedness,
and is open.

% proof of prop:similarity-general-index-zero
\begin{proof}[Sketch]
\emph{(1).} $\lambda_{\min}(\Sigma)I \le \Sigma \le \lambda_{\max}(\Sigma)I$, so by
Proposition~\ref{prop:similarity-general-sandwich}(2) $L_\Sigma$ lies between $L_{tI} = (\sqrt t)_*L_I$ at the
two values of $t$. By \ref{eq:polar-decomposition} of Definition~\ref{def:isotropic-origin-datum},
$L_{tI}(K_A(\varepsilon)) = \overline\sigma(A)\int_\varepsilon^1k_d(rt^{-1/2})\,dr/r$, and $k_d \to c$ at the
origin.

\emph{(2).} The same sandwich for the exponents, with $\psi_{tI}(\omega) - \psi_I(\omega) \to \tfrac c2\log t$ by
Lemma~\ref{lem:isotropic-exponent-log}(2).

\emph{(3).} For the step $\Sigma \le tI$, $t := \lambda_{\max}(\Sigma)$, let $\rho$ be its increment under the
second hypothesis. Under the first, let $\rho$ be the non-Gaussian factor of its increment: the law with
triplet $(0, \nu', 0)$ (C2(iii)), where $\nu = L_{tI} - L_\Sigma$ is the \Levy{} measure of the increment
(C2(ii)) and $\nu'$ its symmetrization. In both cases $\widehat\rho = e^{-(\psi_{tI} - \psi_\Sigma)}$, so
$0 < \widehat\rho \le 1$, and $\widehat\rho(\omega/\varepsilon) \to t^{-c/2}e^{\Delta_\Sigma(\omega/|\omega|)}$ by
Lemma~\ref{lem:isotropic-exponent-log}(2) and (D). On the measure side
$\int\phi(x/\varepsilon)\,\rho(dx) \to \phi(0)\rho(\{0\})$ by dominated convergence, and Fourier inversion
makes the limit of the transforms the constant $\rho(\{0\})$ almost everywhere. For an elliptical part of a
profile bounded at the origin (D) holds by Lemma~\ref{lem:isotropic-exponent-log}(2), and this gives
Proposition~\ref{prop:similarity-orbit-family}(4) again.
\end{proof}

\subsection{The families of the scheme}\label{app:similarity-scheme}

The statements above assume that every Loewner step has an infinitely divisible increment. For the
families of the scheme this is a theorem, which discharges that hypothesis everywhere above.

% shared with blueprint prop:similarity-cascade-families+f33d542fd9b1
\begin{proposition}[similarity cascades: infinitely divisible increments are a theorem of the axioms]\label{prop:similarity-cascade-families}
Let $d \ge 2$, let $G$ be a closed subgroup of $\GL(d)$ containing the similarities $cQ$, $c > 0$,
$Q \in \Orth(d)$, and let $(\mu_{\Sigma,\Sigma'})$ be a $(G, \text{Loewner})$-cascade family of kernels
(Definition~\ref{def:gw-cascade-family}(2)), with kernels from scale zero $\mu_\Sigma := \mu_{0,\Sigma}$.
\begin{enumerate}
\item \emph{(Infinitely divisible kernels and increments.)} Every $\mu_{\Sigma,\Sigma'}$ is infinitely divisible. For
every $\Sigma \in \PD(d)$, $\mu_\Sigma$ is a symmetric infinitely divisible law with generating triplet
$(C(\Sigma), L_\Sigma, 0)$, where $C(\Sigma) \in \PSD(d)$ and $L_\Sigma$ is a symmetric \Levy{} measure;
$(\sqrt c\,Q)_*\mu_\Sigma = \mu_{cQ\Sigma Q^{\tp}}$ for $c > 0$ and $Q \in \Orth(d)$; and every Loewner
step $\Sigma \le \Sigma'$ in $\PD(d)$ has the infinitely divisible increment $\mu_{\Sigma,\Sigma'}$,
$\mu_{\Sigma'} = \mu_\Sigma * \mu_{\Sigma,\Sigma'}$.
\item \emph{(The general part.)} Consequently $(\mu_\Sigma)_{\Sigma \in \PD(d)}$ is a family as in
Proposition~\ref{prop:similarity-general-sandwich}, and the hypothesis ``every Loewner step in $\PD(d)$ has an
infinitely divisible increment'' holds. So Proposition~\ref{prop:similarity-general-sandwich}(1)--(5),
Proposition~\ref{prop:similarity-general-bounded-shape}(1)--(4) and Proposition~\ref{prop:similarity-general-index-zero}(1)--(3)
hold for it under their remaining hypotheses alone, which are:
\begin{itemize}
\item the comparability $L_* \ge \theta L_I$ in Proposition~\ref{prop:similarity-general-bounded-shape}(2);
\item $\psi_* \ge \theta\psi_I$ in (3) there;
\item $L_I \ne 0$ in (4) there;
\item in Proposition~\ref{prop:similarity-general-index-zero}, that $L_I$ is the \Levy{} measure of an isotropic
origin datum $(0, k_d)$ with $k_d(0{+}) < \infty$, with (D) in its clause (3).
\end{itemize}
\item \emph{(Elliptical parts.)} If moreover $\mu_\Sigma = N(0, C(\Sigma)) * \nu_{h(\Sigma)}$ for every
$\Sigma \in \PD(d)$, with $k$, $(\nu_H)$, $h$ and $C$ as in Proposition~\ref{prop:similarity-orbit-family}, then that
proposition's clauses (2) and (3) hold under their remaining hypotheses alone (in (2): $h$ equivariant and
extending continuously to $\PSD(d) \setminus \{0\}$ with values in $\PD(d)$; in (3): $k(0{+}) < \infty$
and $h$ equivariant), with or without a Gaussian part.
\end{enumerate}
\end{proposition}

So for every family of the scheme at (similarities, Loewner), with or without a Gaussian part, the shape
is bounded whenever the rank-one limit is comparable to the isotropic member, and at index $0$ the part
is shape-blind at the origin. Only the infinite divisibility is discharged. The comparability
conditions, $L_I \ne 0$, the origin datum and (D) remain hypotheses. The statements above keep their
hypothesis because they concern families of kernels from scale zero with no increments between other
pairs and no continuity; replacing it by the cascade axioms could only narrow them.

% proof of prop:similarity-cascade-families
\begin{proof}
\emph{(1).} $(G, \text{Loewner})$ is a wedge datum (Definition~\ref{def:gw-cascade-family}), and by
Theorem~\ref{thm:cascade-infinitely-divisible}(3)(a) it has the path condition, so by
Theorem~\ref{thm:cascade-infinitely-divisible}(2) every $\mu_{\Sigma,\Sigma'}$ and every
$\mu_\Sigma = \mu_{0,\Sigma}$ is infinitely divisible. The similarity $-I$ fixes every $\Sigma$, so (K3)
makes $\mu_\Sigma$ symmetric, and the uniqueness C2(ii) makes its drift $0$ and its \Levy{} measure
symmetric. (K3) at $(0, \Sigma)$ with $A = \sqrt c\,Q$ is the covariance, and (K1) at
$0 \le \Sigma \le \Sigma'$ gives the increment.

\emph{(2)} and \emph{(3).} By (1) the family is as in
Proposition~\ref{prop:similarity-general-sandwich}, and every Loewner step has an infinitely divisible
increment, which is the hypothesis of the named clauses of
Propositions~\ref{prop:similarity-general-sandwich}, \ref{prop:similarity-general-bounded-shape},
\ref{prop:similarity-general-index-zero} and~\ref{prop:similarity-orbit-family}.
\end{proof}
